% ECONOMICS_PROBLEM: {"id": "O47-596", "title": "Productivity slowdown and intangible output", "jel_code": "O47", "source_id": "OP-0017"}
% !TeX program = xelatex
\documentclass[11pt,a4paper]{article}
\usepackage[margin=23mm]{geometry}
\usepackage{fontspec}
\setmainfont{Latin Modern Roman}
\usepackage{amsmath,amssymb,mathtools}
\usepackage{xcolor,enumitem,booktabs,longtable,array,xurl,hyperref}
\definecolor{muted}{HTML}{53636C}
\hypersetup{colorlinks=true,urlcolor=blue,linkcolor=blue}
\setlength{\parindent}{0pt}
\setlength{\parskip}{5pt}
\newcommand{\R}{\mathbb R}
\newcommand{\E}{\mathbb E}
\newcommand{\N}{\mathbb N}
\newcommand{\ind}{\mathbf 1}
\newcommand{\Var}{\operatorname{Var}}
\newcommand{\Cov}{\operatorname{Cov}}
\newcommand{\supp}{\operatorname{supp}}
\DeclareMathOperator*{\argmin}{arg\,min}
\DeclareMathOperator*{\argmax}{arg\,max}
\newcommand{\entrymeta}[3]{{\sffamily\small\color{muted}\textbf{\texttt{#1}}\quad|\quad #2\par #3\par}}
\newcommand{\Problemref}[1]{\texttt{#1}}
\newcommand{\JELref}[2]{\texttt{#2} (JEL \texttt{#1})}
\begin{document}
\section*{Productivity slowdown and intangible output}
\textbf{Problem ID:} \texttt{O47-596}\quad \textbf{JEL:} \texttt{O47}\par
% BEGIN ECONOMICS STATEMENT
\label{problem:OP-0017}
\label{atlas:prob:A7}

\noindent\textbf{Original atlas ID:} A7. \textbf{Formulation:} editorial specification of a research family.

\subsection*{Definitions and assumptions}

Fix a country, sample, hours series $H_t>0$, and a conventional real-output index $Y_t^M>0$. Define measured labor-productivity growth $g_t^M=\Delta\log(Y_t^M/H_t)$. An augmented production-accounting model $A$ supplies a consistent output index $Y_t^A$ and tangible/intangible capital services, with capitalization, depreciation, and price-deflator conventions stated explicitly. Let $g_t^A=\Delta\log(Y_t^A/H_t)$. For prespecified periods $P_0,P_1$, bars denote their average growth, and let the measured slowdown be $D^M=\bar g_{P_0}^M-\bar g_{P_1}^M\ne0$. Separately, let $V_i$ be household $i$'s indirect utility and $m_i$ its income. Benefit from access to a free good is compensating willingness to pay $c_i$, when well defined under strict monetary monotonicity, satisfying $V_i(\text{access},m_i-c_i)=V_i(\text{no access},m_i)$ under a declared preference and income model.

\subsection*{Formal research question}

Identify the contribution of accounting changes to the measured slowdown:
\[
\rho_A=\frac{D^M-(\bar g_{P_0}^A-\bar g_{P_1}^A)}{D^M}.
\]
Characterize the set of $\rho_A$ over augmented-accounting models consistent with observed prices, quantities, investment, and validated quality adjustments. Determine which missing intangible investment or digital-quality changes explain a material portion of the slowdown, with uncertainty for both historical periods. Estimate free-good benefits in a separate welfare account, and identify assumptions under which a welfare growth index can be compared with production growth. Report accounting reclassification separately from an economically estimated change in productive technology or business efficiency.

\subsection*{Interpretation and scope}

The ratio $\rho_A$ need not lie between zero and one: a correction can amplify the slowdown or reverse it. Treating it automatically as a bounded causal fraction would be incorrect. GDP measures production according to specified accounting conventions, while consumer surplus from free services measures a distinct welfare object. Adding survey willingness to pay directly to conventional real GDP can mix prices, quantities, and benefits without a consistent index-number derivation. Solow residuals also capture measurement error, utilization, market power, and omitted inputs unless these are restricted. Meaningful comparisons require consistent treatment of capitalizing intangibles in both output and capital services, rather than inserting extra assets into only one side of the growth account.

\subsection*{Source audit and status}

Solow (1957), Corrado--Hulten--Sichel (2009), and Syverson (2017) are verified. Syverson provides evidence challenging particular mismeasurement accounts, so the atlas cannot imply that unmeasured digital value is already established as the main explanation. Original link [19] is now verified as Brynjolfsson--Collis--Diewert--Eggers--Fox (2025), \emph{GDP-B: Accounting for the Value of New and Free Goods}, a published welfare-accounting proposal listed as [S4]. It motivates the separate benefits target rather than a mechanical GDP addition. The slowdown ratio and admissible-accounting model set above are editorial; unresolved measurement contributions must be investigated for a specified country and period.

\subsection*{References}

\begin{enumerate}

\item[S1] Robert M. Solow (1957). \emph{Technical Change and the Aggregate Production Function}. Review of Economics and Statistics 39(3), 312--320. \url{https://www.jstor.org/stable/1926047}. \textit{Locator:} article; growth-accounting derivation. \textit{Establishes:} A foundational productivity-accounting decomposition. \textit{Does not establish:} that the residual measures pure technology without additional assumptions.

\item[S2] Carol Corrado, Charles Hulten, and Daniel Sichel (2009). \emph{Intangible Capital and U.S. Economic Growth}. Review of Income and Wealth 55(3), 661--685. \url{https://onlinelibrary.wiley.com/doi/10.1111/j.1475-4991.2009.00343.x}. \textit{Locator:} abstract; augmented capital accounting. \textit{Establishes:} Material effects of including specified intangible assets in growth accounts. \textit{Does not establish:} that every digital consumer benefit is market output.

\item[S3] Chad Syverson (2017). \emph{Challenges to Mismeasurement Explanations for the US Productivity Slowdown}. Journal of Economic Perspectives 31(2), 165--186. \url{https://pubs.aeaweb.org/doi/10.1257/jep.31.2.165}. \textit{Locator:} abstract; four empirical challenges. \textit{Establishes:} Evidence against particular large mismeasurement explanations of the slowdown. \textit{Does not establish:} zero measurement error or a complete post-2017 explanation.

\item[S4] Erik Brynjolfsson, Avinash Collis, W. Erwin Diewert, Felix Eggers, and Kevin J. Fox (2025). \emph{GDP-B: Accounting for the Value of New and Free Goods}. American Economic Journal: Macroeconomics 17(4), 312--344. \url{https://www.aeaweb.org/articles?id=10.1257/mac.20210319}. \textit{Locator:} abstract; welfare-accounting framework. \textit{Establishes:} A benefit metric and experimental valuations for new/free goods. \textit{Does not establish:} that welfare benefits can be added mechanically to conventional real GDP.

\end{enumerate}

\subsection*{Common conventions: Source mathematical conventions}

Unless an entry states otherwise, agent and firm sets are finite. State and action spaces are measurable, strategies and outcome maps are measurable, probability laws are specified, and expectations exist for the targets under discussion. Infinite-horizon discount factors lie in $(0,1)$ and discounted objectives are finite. Norms, tolerances and complexity input sizes must be specified for computational comparisons. Constants or primitives that appear locally are inputs, not empirical facts asserted by this catalogue.

\textbf{Identification.} A statistical model $m\in\mathcal M$ implies an observable law $P_m$ and target $\theta(m)$. At law $P$, its identified set is
\[
\Theta(P;\mathcal M)=\{\theta(m):m\in\mathcal M,\ P_m=P\}.
\]
Point identification holds when this set is a singleton, or equivalently when $P_{m_1}=P_{m_2}$ implies $\theta(m_1)=\theta(m_2)$ for all compatible models. Sharp bounds mean bounds attained, or approached when the boundary is not attained, by compatible models. Writing this set is a definition, not a solution: the substantive task is to establish informative restrictions and derive or estimate the set. An unrestricted-confounding model may give only trivial bounds.

\textbf{Inference.} Identification, estimability and valid uncertainty quantification are separate questions. Fix $\alpha\in(0,1)$. A confidence procedure $C_n$ over a declared law class $\mathcal P$ has asymptotic uniform coverage at level $1-\alpha$ when
\[
\liminf_{n\to\infty}\inf_{P\in\mathcal P}\Pr_P\{\theta(P)\in C_n\}\geq1-\alpha.
\]
For set-identified targets the coverage event must be defined separately, for example coverage of the full identified set. If an entry uses identification-indexed models instead of uniquely indexed laws, the same coverage convention is applied over those models. Pointwise limits do not establish uniform validity, and asymptotic validity does not guarantee finite-sample precision. Sample sizes, dependence structures and weak-identification sequences are part of each application.

\textbf{Counterfactuals.} $Y_i(a)$ is a potential outcome under a specified intervention, including its timing, implementation and versions. It is not $Y_i$ conditional on voluntarily choosing $a$. A full intervention vector permits interference; shorthand scalar treatment notation does not silently impose no interference. Assignment, selection, attrition, anticipation and measurement must be specified. A claim about an instrument's local effect is not automatically a population policy effect.

\textbf{Economic equilibrium.} A counterfactual model includes best responses, constraints, feasibility, market clearing, state transitions and expectations consistent with the stated information. When several equilibria exist, either a declared selector is an input or the target is set-valued. Comparisons across policy regimes must use the stated selection convention. Reduced-form relationships are not presumed invariant after an equilibrium-changing intervention.

\textbf{Optimization and welfare.} Optimization is over the stated feasible set. If attainment is not established, $\sup$ or $\inf$ and $\varepsilon$-optimal choices replace an asserted maximizer or minimizer. For example, with value $V=\sup_{a\in\mathcal A}W(a)<\infty$, an $\varepsilon$-optimal set is $\{a\in\mathcal A:W(a)\geq V-\varepsilon\}$ for $\varepsilon>0$. Benefits, costs, risk and health or environmental outcomes can be aggregated only using declared units and conversion weights. Interpersonal utility weights, the treatment of fiscal transfers and foreign profits, discounting, and the behavioral welfare criterion are explicit inputs. A comparative-static derivative additionally requires existence and appropriate differentiability of the selected counterfactual.

\textbf{Sources.} Local labels [S1], [S2], and so forth restart in every question. Locators identify the relevant abstract, model, section or result at the level actually checked; a bibliographic or abstract-level verification is not represented as inspection of an entire proof. Working papers are distinguished from later publications and changing versions. Source summaries are paraphrased. All URL checks and editorial formulations refer to the audit date stated above.
% END ECONOMICS STATEMENT
\end{document}
